\uuid{oCr3}
\titre{Limite en l'infini}
\auteur{bodin ; hanani ; mzari}
\organisation{exo7}
\licence{CC BY-NC-SA 4.0}
\datecreate{2020-01-20}
\niveau{L1}
\module{Analyse}
\chapitre{Continuité, limite et étude de fonctions réelles}
\sousChapitre{Limite de fonctions}
\difficulte{1}
\theme{Limites des fonctions réelles, limite en l'infini}
\exo7qcmid{lille-1.limites.008}

\contenu{
\begin{question}
\qclassification{123.03}
\qtags{motcle=limite en l'infini}
Soit $f(x)= (x^5-x^3+1)e^{-x}$. Quelles sont les assertions vraies ?
\begin{answers}
    \good{$\lim_{x\to +\infty} f(x)=0$}

    \bad{$\lim_{x\to +\infty} f(x)=+\infty$}

    \good{$\lim_{x\to -\infty} f(x)=-\infty$}

    \bad{$\lim_{x\to -\infty} f(x)=+\infty$}
\end{answers}
\begin{explanations}
Si $\alpha$ et $ \beta$ sont des réels $>0$,
alors  en $+\infty$, on a :
$ x^{\alpha}\ll  e^{\beta x}$, où la notation $f\ll g$ signifie : $\lim_{x \to + \infty} \frac{f(x)}{g(x)}=0 $.
\end{explanations}
\end{question}
}
